\newpage
\section{SWIRL-based VM security level calculation}\label{sec:swirl_security_level}

The SWIRL proof system is used by OpenVM2. It is a matrix-stacking proof system
that reduces trace evaluations to a stacked polynomial evaluation, which is then
opened via the WHIR PCS (\cref{sec:whir_security_level}).

This companion summarizes the bounds implemented by soundcalc. For full details, see:
\begin{itemize}
    \item SWIRL paper: \url{https://openvm.dev/swirl.pdf}
    \item Canonical OpenVM2 soundness calculator:
          \url{https://github.com/openvm-org/stark-backend/blob/develop-v2/crates/stark-backend/src/soundness.rs}
\end{itemize}

The implementation follows the canonical calculator round by round.\footnote{See
\texttt{soundcalc/circuits/swirl/calculator.py:calculate\_swirl\_soundness()}.}
Let \(L_{\mathrm{PCS}}\) be the list-size bound for the initial WHIR proximity
regime. For unique decoding, \(\log_2 L_{\mathrm{PCS}}=0\). For list decoding,
the BCHKS25 bound used for WHIR determines \(\log_2 L_{\mathrm{PCS}}\).

\subsection{SWIRL component bounds}

LogUp soundness includes the PCS list-size union bound and effective proof-of-work
bits \(g_{\mathrm{eff}}\), where \(g_{\mathrm{eff}}\) accounts for the small
\texttt{sample\_bits} bias over the BabyBear base field:\footnote{See
\texttt{soundcalc/circuits/swirl/calculator.py:SWIRLLogUpSecurityParameters.get\_soundness\_bits()}
and \texttt{soundcalc/circuits/swirl/calculator.py:\_effective\_pow\_bits()}.}
\[
\lambda_{\mathrm{LogUp}}
= \log_2 |\mathbb{F}_{\mathrm{ext}}|
  - \log_2(2 I_{\max})
  - \log_2 M_{\max}
  - \log_2 L_{\mathrm{PCS}}
  + g_{\mathrm{eff}}.
\]

The ordinary GKR checks use the same worst per-round bounds as the backend:
\[
\lambda_{\mathrm{GKR,sumcheck}} =
\log_2 |\mathbb{F}_{\mathrm{ext}}| - \log_2 3,
\qquad
\lambda_{\mathrm{GKR,batching}} =
\log_2 |\mathbb{F}_{\mathrm{ext}}|.
\]

ZeroCheck after the fused input-layer boundary has degree
\[
d_{\mathrm{zc}} =
\max\{(d_{\max}+1)(2^{\ell_{\mathrm{skip}}}-1),\,d_{\max}+1\},
\]
so\footnote{See
\texttt{soundcalc/circuits/swirl/calculator.py:calculate\_zerocheck\_sumcheck\_soundness()}.}
\[
\lambda_{\mathrm{ZeroCheck}}
= \log_2 |\mathbb{F}_{\mathrm{ext}}| - \log_2 d_{\mathrm{zc}}
  - \log_2 L_{\mathrm{PCS}}.
\]

The fused boundary and batch-sumcheck batching term uses\footnote{See
\texttt{soundcalc/circuits/swirl/calculator.py:calculate\_constraint\_batching\_soundness()}.}
\[
d_{\mathrm{fused}}
=
\max\{3,n_{\mathrm{extra}}\}
+(2^{\ell_{\mathrm{skip}}}-1)
+(C_{\max}-1),
\qquad
d_{\mathrm{batch}} = 3A-1,
\]
where \(C_{\max}\) is the maximum number of constraints in one AIR and \(A\) is
the number of AIRs. The security is
\[
\lambda_{\mathrm{fused}}
=
\min\{
\log_2 |\mathbb{F}_{\mathrm{ext}}|-\log_2 d_{\mathrm{fused}},
\log_2 |\mathbb{F}_{\mathrm{ext}}|-\log_2 d_{\mathrm{batch}}
\}
-\log_2 L_{\mathrm{PCS}}.
\]

Stacked reduction uses the minimum of trace-column batching, the univariate
round, and multilinear rounds, again with the PCS list-size union bound:
\[
\lambda_{\mathrm{stack}}
=
\min\{
\log_2 |\mathbb{F}_{\mathrm{ext}}|-\log_2(2N_{\mathrm{cols}}),
\log_2 |\mathbb{F}_{\mathrm{ext}}|-\log_2(2(2^{\ell_{\mathrm{skip}}}-1)),
\log_2 |\mathbb{F}_{\mathrm{ext}}|-1
\}
-\log_2 L_{\mathrm{PCS}}.
\]

\subsection{SWIRL WHIR bounds}

SWIRL uses the OpenVM backend WHIR soundness calculation rather than the older
generic WHIR routine.\footnote{See
\texttt{soundcalc/circuits/swirl/calculator.py:\_calculate\_whir\_soundness()}.}
Each folding sub-round combines the WHIR sumcheck error with the proximity-gap
error by summing the two errors before converting to bits:
\[
\lambda_{\mathrm{fold}}
= -\log_2(2^{-\lambda_{\mathrm{sumcheck}}}
          +2^{-\lambda_{\mathrm{gap}}}).
\]
Each shift or final round similarly combines query sampling with in-domain
\(\gamma\)-batching:
\[
\lambda_{\mathrm{shift}}
= -\log_2(2^{-\lambda_{\mathrm{query}}}
          +2^{-\lambda_{\gamma}}).
\]
The WHIR contribution is the minimum of \(\mu\)-batching, all folding
round-by-round bounds, all shift/final bounds, and non-final OOD bounds.

\subsection{SWIRL proof size}

For OpenVM2, soundcalc estimates SWIRL proof size by counting the backend proof
codec fields, not by applying the generic WHIR proof-size model alone.\footnote{See
\texttt{soundcalc/circuits/swirl/proof\_size.py:get\_swirl\_proof\_size\_breakdown\_bits()}.}
The total proof size is decomposed as
\[
\bproof =
b_{\mathrm{pre}} + b_{\mathrm{GKR}} + b_{\mathrm{batch}}
+ b_{\mathrm{stack}} + b_{\mathrm{WHIR}}.
\]
The terms correspond to the top-level fields of
\texttt{stark-backend/src/proof.rs}: the proof preamble, LogUp GKR proof,
batch-constraint proof, stacking proof, and backend WHIR proof. Base-field
elements are counted by their encoded byte length; for BabyBear this is
\(32\) bits even though the field modulus has \(31\) bits.

The OpenVM2 configuration uses proof-size bounds that assume all optional AIRs
are present. If no explicit \texttt{proof\_size\_*} value is provided for a
circuit, the proof-size calculator uses the matching \texttt{soundness\_*}
bound where available.
